В формальной модели предполагается, что ориентированный ациклический граф
$G=(V,E)$ известен до начала планирования и остаётся неизменным в ходе
исполнения. Это допущение приемлемо для заранее специфицированного набора работ,
но ограничивает применение модели к исследовательским задачам, архитектурным
разведкам и другим R\&D-процессам, в которых новые задачи и зависимости
обнаруживаются по мере получения результатов. В таких процессах оценки
$L_4(P)$ и $B_4$ относятся только к текущему снимку графа $G_t$; однажды
вычисленный критический путь нельзя интерпретировать как постоянную
характеристику всего проекта.

Все задачи считаются доступными в момент начала планирования, их фазы ---
непрерываемыми, а ограниченными ресурсами --- только $P$ агентных потоков и один
разработчик. Release dates, отдельные мощности CI/API и полностью ручные задачи
требуют другой постановки; M4 описывает только заранее выбранный
AI-поддержанный набор.

Модель также не описывает механизм раскрытия новых вершин, стоимость
перепланирования и выбор способа дальнейшей декомпозиции. Дробление задачи может
выявить только содержательно допустимый параллелизм: оно не устраняет реальные
отношения предшествования и обычно увеличивает затраты на постановку, проверку
и интеграцию. Поэтому приведённые результаты позволяют оценивать уже заданный
граф или его текущий снимок, но сами по себе не определяют оптимальную политику
построения динамического графа.

Абсолютный прогноз времени требует калибровки $X$, $k_i$, $h_i$, $C(P)$ и
$\gamma(P)$ для конкретных инструментов и режима работы; при их смене такая
калибровка может быстро устаревать. Для сравнительного выбора между двумя
вариантами требования слабее: объёмы и длительности можно задать в условных
единицах одной базовой задачи, если обе конфигурации оценены в общей шкале.
Предложение~\ref{prop:invariance} гарантирует устойчивость ранжирования только
к одному общему множителю, одинаково применённому ко всем agent- и human-фазам
сравниваемых вариантов. Неодинаковые ошибки по задачам или режимам, искажение
разбиения между $a_i$ и $h_i$, внутреннего фазового профиля, а также разные
значения $C(P)$ и $\gamma(P)$ способны изменить ресурсные пути и порядок
вариантов.
